\item \textbf{ACTIVIDAD: Los datos refractivos de Claudio Ptolomeo.}
Su grupo ha sido asignado para aprender sobre Claudio Ptolomeo ($100-170$) y analizar sus primeros resultados sobre la refracción de la luz. En su libro \textit{Óptica}, Ptolomeo incluye una tabla de ángulos de incidencia y ángulos refractados para la luz que entra al agua. Este es el conjunto de datos más antiguos que sobreviven de este tipo. Los datos de Ptolomeo se muestran en la siguiente tabla:

\begin{center}
\begin{tabular}{|c|c|}
\hline
\textbf{Ángulo de incidencia ($\theta_1$)} & \textbf{Ángulo de refracción ($\theta_2$)} \\ \hline
$10.0^\circ$ & $8.0^\circ$ \\ \hline
$20.0^\circ$ & $15.5^\circ$ \\ \hline
$30.0^\circ$ & $22.5^\circ$ \\ \hline
$40.0^\circ$ & $29.0^\circ$ \\ \hline
$50.0^\circ$ & $35.0^\circ$ \\ \hline
$60.0^\circ$ & $40.5^\circ$ \\ \hline
$70.0^\circ$ & $45.5^\circ$ \\ \hline
$80.0^\circ$ & $50.0^\circ$ \\ \hline
\end{tabular}
\end{center}

¿Cuál es el índice de refracción del agua según los datos de Ptolomeo?